\documentclass[11pt,a4paper]{article}
\usepackage[utf8]{inputenc}
\usepackage[T1]{fontenc}
\usepackage[portuguese]{babel}
\usepackage{lmodern}
\usepackage[margin=2.3cm]{geometry}
\usepackage{tabularx,longtable,booktabs,array,enumitem}
\usepackage[table]{xcolor}
\usepackage{tcolorbox}
\usepackage{hyperref}
\definecolor{roamly}{HTML}{0E6E68}
\definecolor{alerta}{HTML}{B3261E}
\definecolor{formal}{HTML}{2E7D32}
\hypersetup{colorlinks=true,linkcolor=roamly,urlcolor=roamly}
\newtcolorbox{aviso}[1][]{colback=alerta!6,colframe=alerta,arc=2pt,boxrule=0.8pt,fonttitle=\bfseries,title=#1}
\newcolumntype{L}{>{\raggedright\arraybackslash}X}
\renewcommand{\arraystretch}{1.3}
\setlist{nosep,leftmargin=*}
\newcommand{\bolt}{\textcolor{alerta}{[BD]}}

\title{\textbf{Roamly}\\[2pt]\large Entrevista a jovem viajante: U0\\\normalsize Registo da entrevista}
\author{}\date{}
\begin{document}
\maketitle

\begin{aviso}[Caso 2: Jovem viajante com orçamento limitado]
Tema: \textbf{planear uma viagem em grupo dentro do orçamento}.
Anónimo: sem nomes de pessoas, instituições ou cidade de origem; os amigos são
«Amigo~A» e «Amigo~B». O código \texttt{U0} é reservado a este caso.
\end{aviso}

%==========================================================
\section{Perfil e situação da entrevista}
\begin{tabularx}{\textwidth}{@{}p{4.3cm}L@{}}
\toprule
Código & U0 \\
Papel (Cap.~4) & Organizador(a) do grupo, um(a) caçador(a) de promoções \\
Perfil & 20 anos, 2.º ano de licenciatura, vive em apartamento partilhado, 
sem aeroporto na cidade; mesada (\(\approx\)250\,€) e explicações pontuais 
(\(\approx\)120\,€); poupança para viajar \(\approx\)380\,€; cartão de débito e 
MB Way, sem crédito \\
Viagem em causa & 4 noites em Cracóvia com 2 amigos; limite \textbf{300\,€ por pessoa},
tudo incluído \\
Tipo de tarefa & Longa e intermitente: \emph{work walkthrough} da viagem já planeada; 
U0 mostra e refaz os passos no portátil e no telemóvel \\
Local e duração & No quarto de U0, com portátil e telemóvel; \(\approx\)1\,h45, 
dos quais \(\approx\)1\,h30 de relato guiado \\
Registo & Transcrição do conteúdo considerado necessário (com consentimento); sem 
capturas de ecrã \\
\bottomrule
\end{tabularx}

%==========================================================
\section{Siglas}
{\small
\begin{tabularx}{\textwidth}{@{}p{2.6cm}L@{}}
\toprule
\textbf{Sigla} & \textbf{Significado} \\ \midrule
\multicolumn{2}{@{}l}{\textit{Identificadores}} \\
U0 & Código do viajante entrevistado \\
UE01--UE19 & Registos da entrevista ao viajante (observações e relatos) \\
U0-01--U0-21 & Notas da sessão de interpretação do caso U0 \\ \midrule
\multicolumn{2}{@{}l}{\textit{Tipos de registo da entrevista}} \\
O & Observado: U0 mostra ou refaz o passo durante a entrevista \\
R & Relato de caso concreto, sem demonstração \\
F/OP & Fala geral ou opinião, sem caso concreto (fora dos modelos) \\ \midrule
\multicolumn{2}{@{}l}{\textit{Tipos de nota de interpretação}} \\
OBS & Observação \\
INS & \textit{Insight} \\
CUL & Influência cultural \\
Q & Pergunta por esclarecer \\
BD & \textit{Breakdown}: falha ou atrito na tarefa \\
OP & Opinião geral \\ \midrule
\multicolumn{2}{@{}l}{\textit{Outras}} \\
\bolt{} & Marca de \textit{breakdown} nos registos UE \\
IA & Inteligência artificial \\
\bottomrule
\end{tabularx}}

%==========================================================
\section{Conversa convencional (\(\approx\)10\,min)}
Tipo F/OP: contextualiza mas \textbf{não entra nos modelos} sem caso concreto.

\begin{tabularx}{\textwidth}{@{}p{1.6cm}L@{}}
\toprule
\textit{Entrev.} & «Como costumas planear uma viagem?» \\
U0 & «\textcolor{formal}{Para cada parte da viagem comparo em imensos sites 
e escolho o mais barato.}» 
(\textbf{OP}) \\
\textit{Entrev.} & «E no teu orçamento, de onde vem o dinheiro?» \\
U0 & «Mesada, algumas explicações e o que tenho guardado. Guardo para viajar porque é o que mais gosto de fazer.» \\
\textit{Entrev.} & «O que achas das ferramentas que usas?» \\
U0 & «\textcolor{formal}{Funcionam, mas o preço que aparece nunca é o preço final. E tenho tudo espalhado: voos num sítio, quartos noutro, contas numa folha.}» (\textbf{OP}) \\
\bottomrule
\end{tabularx}

\paragraph{Transição.} «Vou pedir-te para me mostrares, passo a passo, como planeaste esta viagem, usando o que tens aí. Eu vou interrompendo para perceber.»

%==========================================================
\section{Entrevista contextual: relato guiado}
Ordem cronológica da viagem. \textbf{UE} = registo da entrevista ao viajante; \textbf{O} = U0 mostra ou refaz o passo; \textbf{R} = relato de caso concreto. Itálico = pergunta do entrevistador.

{\small
\begin{longtable}{@{}p{1.1cm}p{0.7cm}p{13.4cm}@{}}
\toprule
\textbf{ID} & \textbf{Tipo} & \textbf{Registo} \\ \midrule \endhead
UE01 & R & À noite, no sofá, telemóvel. Viu um vídeo de uma rede social «Cracóvia 
por 150\,€» e guardou-o. Criou na hora um grupo de mensagens com o Amigo~A e o 
Amigo~B e enviou o vídeo com «bora?». Foi a primeira mensagem do grupo. \\

UE02 & R & Amigo~A: resposta imediata. Amigo~B: «depende
do preço; era mais fácil perto do fim de semana». \\

UE03 & O & Refaz a pesquisa de terça à noite: abre o comparador de voos com 
origem «aeroporto mais próximo» («não há aeroporto na minha cidade»), destino 
Cracóvia, datas «mês inteiro». \\

UE04 & O & Três combinações de datas. Hesita \(\approx\)20\,s no voo mais 
barato (6\,h50). \emph{O que estás a pensar?} «Como é que chego ao aeroporto às 
cinco da manhã?» Pesquisa comboio + metro (\(\approx\)13\,€ por sentido) e 
autocarro (8\,€, mais 1\,h). Escolhe o voo das 11\,h30. \bolt{} voo mais barato 
exige chegar de madrugada sem transporte. \\

UE05 & O & Mostra os dois alojamentos que considerou: dormitório (14\,€/noite) e 
quarto privado para 3 (24\,€ por pessoa e noite). No grupo, o Amigo~A respondeu 
«dormitório não». Ficou o quarto privado. \\

UE06 & O & Mostra a folha de cálculo: duplicou a da viagem anterior («já tem as 
colunas que uso») e copiou à mão os preços de cada site. Total: \textbf{303\,€} 
por pessoa (acima de 300\,€). \\

UE07 & O & Refaz o passo de pagamento do voo e surge a bagagem de porão (+34\,€): 
«isto ninguém diz». Decidiu levar só mochila pequena. \bolt{} custo escondido só 
no pagamento. \\

UE08 & R & Dois dias sem resposta do grupo. Voltou ao comparador: o voo subiu 19\,€. 
Enviou «malta, subiu». \bolt{} \\

UE09 & R & O Amigo~B só confirmou que podia a partir de uma sexta concreta. Refez a 
pesquisa: total \textbf{322\,€}. \\

UE10 & R & Num intervalo entre aulas, no telemóvel, pediu ao ChatGPT um 
«roteiro de 4 dias em Cracóvia, barato». Copiou 4 itens para as notas: «não sei se 
isto é verdade, vou ver no mapa». Confirmou 2 e descartou 1. \\

UE11 & O & Abre o vídeo guardado e mostra o sítio de comida de rua recomendado 
pelo criador de conteúdo, que acrescentou à lista do mapa. \\

UE12 & O & Mostra a linha riscada da visita à mina de sal (28\,€): «isto é para 
turistas». Total \textbf{294\,€}. \\

UE13 & R & No domingo anterior, a mãe ao telefone: «não gastes o que não tens; 
manda-me onde vais dormir». U0 enviou-lhe a ligação do alojamento. \\

UE14 & R & Apareceu o aviso «só restam 2 quartos!». Fechou o separador e abriu 
outro site para comparar: «isto é para me pressionar». \\

UE15 & R & Reservou o voo dos \textbf{três} com o seu cartão de débito (adiantou 
o dinheiro). Alojamento a pagar no local, com pequena caução. \\

UE16 & R & Pediu 98\,€ por MB Way a cada amigo. Amigo~A pagou em minutos. Amigo~B: 
«só no fim do mês». \bolt{} \\

UE17 & R & Partilhou a ligação da folha no grupo; ninguém a abriu. Amigo~A: 
«manda print». Enviou uma captura de ecrã. \bolt{} \\

UE18 & R & Na semana seguinte, no café ao lado da faculdade, o grupo decidiu as
atividades na aplicação GetYourGuide, com um telemóvel passado de mão em mão. Ninguém escreveu; 
depois U0 registou nas notas «o que combinámos». \\

UE19 & O & No fim do relato: «uma colega pagou 360\,€ por quase o mesmo; eu 
consegui por 294 com quarto privado.» Mostra orgulho. \\
\bottomrule
\end{longtable}}

%==========================================================
\section{Fecho: resumo validado com U0}
\begin{itemize}
  \item «O que mais te custa não é procurar, é juntar tudo num total e 
  fazer o grupo decidir a tempo.» U0 corrige: «e ter de adiantar eu o dinheiro.»
  \item «O voo e o alojamento foram escolhas de preço, mas sempre a pensar nos 
  amigos.» U0 confirma.
  \item «O assistente de IA ajudou a ter ideias, mas verificas no mapa.» 
  U0 confirma: «não confio nos preços dele.»
\end{itemize}

%==========================================================
\section{Notas da sessão de interpretação}
Tipos: \textbf{OBS} observação; \textbf{INS} insight; \textbf{CUL} influência cultural; \textbf{Q} pergunta por esclarecer; \textbf{BD} breakdown; \textbf{OP} opinião geral.

{\small
\begin{longtable}{@{}p{1.5cm}p{1.1cm}p{11.8cm}@{}}
\toprule
\textbf{ID} & \textbf{Tipo} & \textbf{Nota} \\ \midrule \endhead
U0-01 & OBS & O gatilho foi um vídeo numa rede social, não uma pesquisa (UE01). \\
U0-02 & OBS & O grupo nasce com a primeira mensagem; a decisão acontece nele (UE01). \\
U0-03 & CUL & O Amigo~B condiciona as datas: «mais fácil perto do fim de semana» (UE02, UE09). \\
U0-04 & OBS & Parte da folha da viagem anterior em vez de uma folha em branco (UE06). \\
U0-05 & INS & Pesquisa por orçamento e datas, não por destino: «mês inteiro» (UE03). \\
U0-06 & BD & O voo mais barato obriga a chegar ao aeroporto de madrugada, sem transporte (UE04). \\
U0-07 & INS & O custo real está espalhado: comparador, transporte até ao aeroporto, folha de cálculo (UE04, UE06). \\
U0-08 & BD & Custo de bagagem só no pagamento: «isto ninguém diz» (UE07). \\
U0-09 & BD & Os preços mudam enquanto o grupo decide (UE08, UE09). \\
U0-10 & CUL & Pressão para decidir depressa («só restam 2 quartos») sentida como manipulação (UE14). \\
U0-11 & CUL & Orgulho em gastar menos que uma colega (UE19); «para turistas» como critério de corte (UE12). \\
U0-12 & CUL & A mãe: «não gastes o que não tens» e pede o local de alojamento (UE13). \\
U0-13 & CUL & O criador de conteúdo é fonte de ideias (UE01, UE11). \\
U0-14 & OBS & Usa o ChatGPT para ideias e verifica no mapa (UE10). \\
U0-15 & OP & «Não confio nos preços dele.» Sustentado em parte por UE10, onde a desconfiança é sobre o roteiro e não sobre preços. \\
U0-16 & BD & U0 adianta o dinheiro dos três e cobra depois; um amigo paga com atraso (UE15, UE16). \\
U0-17 & OBS & Linhas riscadas na folha para justificar ao grupo o que foi cortado (UE12). \\
U0-18 & BD & O grupo não abre a folha partilhada e pede «print» (UE17). \\
U0-19 & OBS & As atividades são escolhidas ao vivo no café, numa aplicação, com um telemóvel passado de mão em mão; depois U0 escreve o que ficou (UE18). \\
U0-20 & Q & Porque é que o grupo responde devagar? Medo de se comprometer, falta de dinheiro, ou desinteresse? Perguntar em entrevistas a quem «vai atrás». \\
U0-21 & Q & Como chega U0 ao aeroporto no dia da viagem? Seguir numa entrevista futura. \\
\bottomrule
\end{longtable}}

\end{document}
